%% n8c_eval.tex — E6 answer behaviour/abstention/cost, E7 deployment arms,
%% E8 verifier failure, E9 operational boundaries.
%% Every figure is attributed to its artefact; corpora are never mixed.
%% LABEL MAP (2026-11-13 renumbering; labels are historical, titles current):
%%   sec:e5 -> E6 | sec:sysarms -> E7 | sec:e6 -> E8 | sec:e7 -> E9
%%   (in n8a: sec:routing-audit -> E3; in n8b: sec:e3 -> E4, sec:ladder-opt -> E5)
%% 2026-11-16: the 30-cell body table moved to supplement S14; E6 now leads
%% with the marker figure (fig_opt_answer) and the layered gold/marker reading.

\subsection{E6: Answer and trace checks, abstention and cost}
\label{sec:e5}
\label{sec:agentic}\label{sec:abstention}\label{sec:composition}\label{sec:eval-agent}
Two ten-question grids on the arXiv question set compare a platform agent
(P), direct file tools (F)
and a chunk-replica arm (C) that searches an 800-character chunk copy of the
same content through the platform's own vector lane---a same-content,
flat-chunk substrate control for the document-level organisation. The
original grid uses different substrates; the later grid pins the same five
libraries, and its C arm searches those libraries directly. Gold-ID
presence, strict and alias-aware keyword coverage and citation markers are
deterministic checks, not assessments of semantic support. An identifier in
the trace need not have been read; a marker does not establish that a cited
passage exists or entails the answer. Figure~\ref{fig:answer} shows the
per-question marker pattern; Table~S14 gives the per-question records and
Table~S15 the aggregate checks and run-time observations.

In both grids the gold identifier occurred in the answer or trace of all 60
cells. Because an identifier in a trace need not have been read, this
identifier-presence check cannot distinguish document reachability; the
observed separation is instead in document-style citation markers. In the
original grid, marker-bearing
answers were 9/10 for P, 7/10 for F and 2/10 for C
(Figure~\ref{fig:answer}a); strict keyword coverage was 9/10 for each arm
and alias-aware coverage 10/10. Under the common substrate, markers were
9/10 in every arm (Figure~\ref{fig:answer}b), strict coverage 7/10, 7/10
and 8/10, and alias-aware coverage 9/10, 9/10 and 10/10. The original-grid
mean time and cost were 52.6 s and $\$0.711$ per question for P, 22.2 s and
$\$0.487$ for F, and 24.5 s and $\$0.527$ for C; these are single-run
observations. Both grids carry a marker in 18 cells and neither in three
(21/30 agreement); the replication adds a marker in nine cells and removes
none (9--0 discordance, McNemar $p=0.0039$), so the marker gap is
substrate-bound rather than an interface
effect alone. The detailed source-presence audit appears in Tables~S7--S8
of the supplementary inventory and is not an entailment test.

A later ten-question run on the repaired platform, with the same five
libraries, returned the gold identifier and a citation marker in all ten
answers, with full keyword coverage in nine (one question covered one of
two annotated terms). Its mean wall-clock was 131 s per question and its
total cost was 3.10 USD. This is a single current-state observation on a
later stack and model state; it is not a controlled arm and is not pooled
with the grids. Read together, the records support a bounded claim:
document-level organisation is associated with document-style citation
behaviour---the tiered platform arm carried markers in 9/10 and, after
repair, 10/10 answers, against 2/10 for the flat chunk replica in the
original grid---but these are lexical proxies; no semantic-quality
superiority is claimed, and the latest run's keyword coverage (9/10) matches
F and trails the historical P arm.

\textbf{Answer-level case study.} A separate recorded demonstration parsed
and ingested a DeepSeek-R1 paper, then used two agent harnesses on the same
four questions. The prompt supplied the library, paper and arXiv identifier
and prescribed the retrieval, judging and reading steps; it did not test
autonomous collection selection. The case compendium reports eight responses
checked against the source on answer content, document provenance, quoted
material and fabricated facts; all eight passed that case review (Table
S16). On the question asking for electricity cost, both responses said the
paper did not provide an electricity figure and distinguished the reported
$\$294$K GPU-rental-based training-cost estimate. This is an observed
refusal to infer an absent value, not a general abstention guarantee. Three
further industrial-corpus probes returned not-found; the retrieval
threshold retained candidates in related absent probes (Fig.~S3).

\begin{figure*}[!tb]
\centering
\includegraphics[width=\textwidth]{fig_opt_answer}
\caption{Answer traceability on two ten-question grids. (a) Original grid:
rows P, F and C, columns BQ01--BQ10; a filled square marks a document-style
marker and the row-end count is the number of marker-bearing answers
(9/10, 7/10, 2/10). (b) Common-substrate replication on the same five
libraries: markers are 9/10 in every arm. Gold-ID presence is 10/10 in every
cell of both grids; strict keyword coverage is 9/9/9 in (a) and 7/7/8 in
(b); alias-aware coverage is 10/10/10 and 9/9/10. Markers are lexical
proxies and do not verify source resolution or entailment; each arm is one
run, and the original grid uses different substrates. The replication
changes the substrate and raises C and F markers, so the marker gap is not
an interface effect alone.}
\label{fig:answer}
\end{figure*}
\begin{figure*}[!tb]
\centering
\includegraphics[width=\textwidth]{fig_opt_verifier}
\caption{Verifier scores and same-pool ordering. (a) Recorded Laya scores
for 191 candidates from sixteen hybrid runs: 17 in-pool golds and 174
non-golds; all recorded values exceed the 0.5 threshold; one additional Q14
fusion candidate has no recorded score. Horizontal jitter is for legibility;
medians are 0.88 and 0.86; the Mann--Whitney $p=0.26$ is exploratory and
ignores clustering by query. (b) The same sixteen pools replayed with fusion
order and with verifier order; filled squares and circles mark a gold in the
top five, open marks a miss; Hit@5 is 15/16 for fusion order and 7/16 for
verifier order. The score archive and the ordering archive are separate
records. The verifier is not a ranker; these records do not show that
threshold filtering is ineffective or that another verifier would behave
the same way.}
\label{fig:verifier}
\end{figure*}

\subsection{E7: Deployment-layer system arms and honesty contracts}
\label{sec:sysarms}
Three questions crossed pinned, unpinned and pinned hybrid-directed API
arms. The pinned arms returned the two in-corpus gold identifiers, and all
three reported not-found for the out-of-corpus probe. In one unpinned
setting, a 131-document, 74.9 KB catalogue exceeded the tool-output budget;
a truncated preview led to false refusal or a hang. After pagination and
grouped views were added, selected reruns completed, including the formerly
failing query. These observations do not establish general transport
reliability; they show that the management-to-answer path can fail in
transport and that the failure is exposed by the workflow rather than
hidden; the silent index-skip fault in E9 motivates surfaced
degradation markers.
An eighteen-arm honesty battery observed no fabricated citations in its
final responses; interrupted arms were rerun and each arm has one final
observation. The deployment checks are operational boundaries, not a
retrieval comparison.

\subsection{E8: When the verifier fails}
\label{sec:e6}
\label{sec:inflation}\label{sec:eval-verify}
The 191 recorded candidate scores from 16 hybrid runs included 17 in-pool
golds and 174 non-golds. All recorded scores exceeded the 0.5 threshold; one
additional Q14 candidate had no recorded score. Gold and non-gold medians
were 0.88 and 0.86 (Figure~\ref{fig:verifier}a). The candidate-level
Mann--Whitney test ($p=0.26$) is exploratory and ignores clustering by
query. The score archive and the ordering archive are separate records.
Re-ranking these scores reduced Hit@5 to 7/16, compared with 15/16 when
recorded survivors retained fusion order (Figure~\ref{fig:verifier}b). A Q2
example appears in Fig.~S5: the gold root-cause report scored 0.846 against
five telemetry reports at 0.887--0.926 and was demoted from fusion rank 1 to
rank 8, outside the read window. A replay of the same pairs with an NLI
checkpoint produced Hit@5 0.3125; a span-level replay produced 0.1875. A
relative-margin replay kept 16 of 17 golds at $m=0.15$ and all 17 at
$m=0.20$, while the mean candidate count stayed near twelve; these replays
neither isolate truncation nor establish the cause of verifier failure. The
arXiv corpus did not show the same top-five collapse; domain, language and
corpus size changed together. No answer-level test validates a higher
filtering threshold; Table~S6 of the supplementary inventory sweeps the
recorded scores and Fig.~S4 replays relative-margin variants of the same
pools. The result is a role warning: the verifier can filter, but ordering
by its score lost the gold in this pool. Filtering itself was never
activated, because every recorded score passed the threshold.

\subsection{E9: Operational boundaries: stability, concurrency, gate and
scope ablations}
\label{sec:e7}
A 27-cell stability campaign completed without recorded answer errors across
three tasks, three lanes and three rounds. On one instance, a concurrency
sweep of 90/180/270 requests at 5/15/30 clients raised median latency from
5.2 to 1577 ms; 111 of 270 requests timed out at 30 clients. A later cache
repair lowered selected 30-client read-ladder median latency to 65.8 ms,
but the full sweep was not repeated. These are single-instance measurements;
raw concurrent-write counts and crash-recovery evidence were unavailable. A
post-acceptance inspection also exposed a coverage fault mode in the
cross-library vector channel: collections whose on-disk index segments are
unreadable are silently skipped inside the per-collection loop, and the
response still reports success without any degraded marker. Logs identified
eight affected collections in one inspection, including a freshly ingested
library. For that library, a scoped probe found the document while the
global search missed it; forced re-indexing restored the global hit. A
scoped probe on another affected library also returned no results, so the
failure cannot be assumed to be confined to the global channel. Logs surface
the underlying read error, but the response does not; until the repair
ships, coverage claims over the global channel should be read against this
fault mode.

A reconstructed fixed-candidate gate comparison used ten questions, a
12-candidate cap and a 4,000-character pack. All three arms packed the gold
first for the eight gold-answerable questions; for nine of ten questions the
filtered and unfiltered document sets matched. LLM grades still varied under
identical packs (including $0\to5$ and $3\to0$), so these observations do
not establish an answer-quality benefit from gating. The two non-gold
questions, an unanswerable clinical question and an out-of-corpus
particle-physics question, were refused in every arm; the judge-survivor
lists were empty in 2/4 A0 and 0/4 A0h out-of-domain routing probes. Both
A0 empty cases lacked a top candidate; these are retrieval-stage outcomes,
not verified gate rejections or answer-level refusal rates.

The scope diagnostic used ten arXiv questions, eight of which had gold
documents in the evaluation corpus; the other two were excluded from Hit@5.
With the gold library pinned, Hit@5 was 1.0 on the eight graded questions;
pinning a non-gold library yielded 0.0. The unscoped search covered 20 live
libraries, including 15 outside the five-library evaluation corpus, and is
not a matched-scope comparison. The catalogue route selected the correct
library first for six of eight answerable questions; file names carried
domain prefixes. These small, single-engine experiments do not establish an
independent effect of content placement or an overall system advantage;
they define the operational envelope in which the management and answer
results should be read.
